% !TeX root = ../../Book3.tex
% 纲要标题：### 12.5 忠实平坦性
% 纲要位置：工作记录/计划纲要.md，第 2194 行。
% 纲要细目（写作范围）：
% * 正合性和消失的反映
% * 忠实平坦基变换下的可检验性质
% * 纤维与下降的基本动机


\section{忠实平坦性}
\label{b3:sec:1205}

平坦基变换保持已有的正合列，却可能把非零模变成零。例如与 \(\mathbb Q\) 张量会使有限挠交换群消失。若希望从变基后的计算推回原来的模，还须保证信息没有这样丢失；忠实平坦性正是这一附加要求。


\subsection{正合性与消失的反映}
\label{b3:subsec:1205:01}

\begin{definition}\label{b3:def:faithfully-flat}
设 \(R\to S\) 为交换环同态。如果 \(S\) 为平坦 \(R\) 模，且对每个 \(R\) 模 \(N\)，从 \(S\otimes_RN=0\) 能推出 \(N=0\)，就称这个同态\term[zhongshi-pingtan]{忠实平坦}[faithfully flat]。
\end{definition}

\begin{theorem}\label{b3:thm:faithfully-flat-criterion}
设 \(R\to S\) 为平坦交换环同态。下列条件等价：
\begin{enumerate}
\item \(R\to S\) 忠实平坦；
\item 对每个极大理想 \(\mathfrak m\subset R\)，有 \(S/\mathfrak mS\ne0\)；
\item \(\operatorname{Spec}S\to\operatorname{Spec}R\) 满射。
\end{enumerate}
特别地，非零局部环之间的平坦局部同态忠实平坦。
\end{theorem}
\begin{proof}
第一项对非零模 \(R/\mathfrak m\) 使用定义，就得到第二项。反过来，设 \(0\ne x\in N\)，则 \(Rx\cong R/I\)，其中 \(I=\operatorname{Ann}_R(x)\) 为真理想。选极大理想 \(\mathfrak m\supseteq I\)，则 \(S/IS\) 满射到非零的 \(S/\mathfrak mS\)，故非零。平坦性使
\(S\otimes_RRx\to S\otimes_RN\) 单射，因而后者非零。这证明第一、二项等价。

若忠实平坦，对每个素理想 \(\mathfrak p\)，非零模 \(\kappa(\mathfrak p)\) 的张量积也非零。非零环 \(S\otimes_R\kappa(\mathfrak p)\) 有素理想，其在 \(S\) 中的收缩恰位于 \(\mathfrak p\) 之上：把纤维写成 \((R\setminus\mathfrak p)^{-1}(S/\mathfrak pS)\)，由局部化素理想对应即可核验。所以谱映射满射。反向，若存在 \(\mathfrak q\subset S\) 收缩为 \(\mathfrak m\)，则 \(\mathfrak mS\subseteq\mathfrak q\ne S\)，于是第二项成立。

局部同态 \((R,\mathfrak m)\to(S,\mathfrak n)\) 满足 \(\mathfrak mS\subseteq\mathfrak n\)，剩余环非零，故最后结论由第二项得到。
\end{proof}

这些判据可对照\cite[\S7, Theorems 7.2--7.3, pp. 47--48]{Matsumura1989CommutativeRingTheory}。检验局部同态时须同时确认平坦和 \(\mathfrak mS\ne S\)；仅仅是局部环之间的一个任意同态并不够。

\subsection{忠实平坦基变换下的性质检测}
\label{b3:subsec:1205:02}

\begin{proposition}\label{b3:prop:faithfully-flat-detection}
设 \(R\to S\) 为忠实平坦交换环同态。下列各项中的模及模同态均在 \(R\) 上，变基后在 \(S\) 上。
\begin{enumerate}
\item 模同态的单射、满射、同构以及模序列的正合性，均可在与 \(S\) 张量后检测。
\item 对每个 \(R\) 模 \(M\)，自然映射 \(M\to S\otimes_RM\)、\(m\mapsto1\otimes m\) 单射；特别地 \(R\to S\) 单射，且对任意理想 \(I\)，有 \(IS\cap R=I\)。
\item \(M\) 在 \(R\) 上平坦，当且仅当 \(S\otimes_RM\) 在 \(S\) 上平坦。
\item 若 \(S\otimes_RM\) 在 \(S\) 上有限生成或有限展示，则 \(M\) 分别在 \(R\) 上有限生成或有限展示。
\end{enumerate}
\end{proposition}
\begin{proof}
平坦性使核、像和余核的短正合列仍然正合。例如
\(S\otimes_R\ker u\cong\ker(1\otimes u)\)，余核也相容，故消失检测给出第一项中的单射、满射与同构。对一个复形 \(L\to M\to N\)，同样把中间同调 \(\ker v/\operatorname{im}u\) 张量，便得正合性检测。若原来还不知道 \(vu=0\)，先对 \(\operatorname{im}(vu)\) 使用同一检测，得到它为零。

若 \(m\ne0\)，平坦性使 \(S\otimes_RRm\) 单射进入 \(S\otimes_RM\)。前者非零，且由 \(1\otimes m\) 作为 \(S\) 模生成，故 \(1\otimes m\ne0\)。这证明第二项的单射；对 \(M=R/I\) 应用它，得到 \(R/I\hookrightarrow S/IS\)，即理想收缩公式。

第三项的正向是\cref{b3:prop:flat-base-change-transitivity}。反向任取单射 \(U\to V\)，先用 \(S\) 的平坦性，继而用 \(S\otimes_RM\) 的 \(S\) 平坦性，得到
\[
S\otimes_R(U\otimes_RM)\longrightarrow
S\otimes_R(V\otimes_RM)
\]
单射。由第一项原映射单射，所以 \(M\) 平坦。

若 \(S\otimes_RM\) 由有限多个张量和生成，收集这些和中出现的有限个 \(m_i\in M\)，令 \(M_0=\sum_iRm_i\)。则 \(S\otimes_R(M/M_0)=0\)，故 \(M=M_0\)，证明有限生成下降。若变基模有限展示，先由刚证结论取 \(R^n\twoheadrightarrow M\)，核为 \(K\)。平坦性给出 \(S\otimes_RK\) 是 \(S^n\to S\otimes_RM\) 的核，由\cref{b3:prop:finite-presentation-kernel}有限生成。再次下降有限生成性，得到 \(K\) 有限生成，故 \(M\) 有限展示。
\end{proof}

结合\cref{b3:thm:finite-presentation-flat-projective}，有限生成投射性也可以忠实平坦地检测。不过“自由”本身不是以上结论的同义词；基变换可能使一个原本只有局部基的模获得全局基。

\subsection{纤维与下降的动机}
\label{b3:subsec:1205:03}

忠实平坦性在任意基变换和复合下保持。复合的证明是依次使用两个正合函子及两个消失检测。对基变换 \(R\to R'\)，\(S'=R'\otimes_RS\) 在 \(R'\) 上平坦；若 \(N\) 为 \(R'\) 模且 \(S'\otimes_{R'}N=0\)，结合律将它识别为 \(S\otimes_RN=0\)，所以 \(N=0\)。

\ProseParagraphBreak
若 \(D(f_1),\ldots,D(f_r)\) 覆盖 \(\operatorname{Spec}R\)，则
\[
R\longrightarrow S=\prod_{i=1}^rR_{f_i}
\]
忠实平坦。有限直积作为模也是有限直和，故平坦；每个素理想至少属于一个 \(D(f_i)\)，局部化的素理想对应说明谱映射满射。单个真局部化一般不忠实，例如 \(\mathbb Z\to\mathbb Z[1/2]\) 遗漏素理想 \((2)\)，并使 \(\mathbb Z/2\mathbb Z\) 消失。

\ProseParagraphBreak
在这个有限开覆盖的例子中，给定变基模相当于在每个开集上给定一个模。要把它们还原为一个全局模，必须说明在重叠的 \(R_{f_if_j}\) 上如何识别，并要求三重重叠上的识别相容。下一节把这个熟悉的拼接问题写成一般忠实平坦环同态上的代数条件。
